\section*{A tensor-prime punctured codimension-two family}

All logarithms are base two and
\(L(t)=t\log_2(1/t)\), with \(L(0)=0\).  For \(n\ge3\), define
\(f_n:\{-1,1\}^n\to\{-1,1\}\) by
\[
 f_n(x)=-1
 \quad\Longleftrightarrow\quad
 x_1=x_2=1\ \hbox{ and }\ x\ne(1,\ldots,1).
 \tag{1}
\]
Thus \(f_n\) is the codimension-two one-point defect with its all-plus
point punctured.  This is a genuine Boolean family in every dimension.

\subsection*{Frozen target and exact implication}

The frozen GU pair is
\[
 G(f)\le4I(f)+\frac34U(f).                              \tag{2}
\]
The independently reviewed row-cut reduction says that (2) follows if
\[
 D_{\ge2}+\frac12V_1
 \le\left(4-\frac2{\ln2}\right)I+\frac34U
      +\frac2{\ln2}Q_{\ge3}                             \tag{3}
\]
for every tensor-prime, positive-residual Boolean function of degree at
least three, where
\[
 Q_{\ge3}=I-M_1-W_2,
 \qquad M_1=\sum_i\min\{p(\varnothing),p(\{i\})\}.
\]
If (3) holds, then \(U\le16I\) and the audited reduction gives
\(H\le34I\).  The purpose of the family (1) is to test the uniformity of
this degree-at-least-three leaf, not to replace its signed left side by an
unrelated entropy.

\subsection*{Complete spectrum}

Put
\[
 \varepsilon=2^{1-n},\qquad r=\varepsilon^2,
 \qquad q_+=(1/2+\varepsilon)^2,
 \qquad q_-=(1/2-\varepsilon)^2.
\]
The unpunctured function is the two-bit one-point defect, extended by dummy
variables.  Puncturing adds twice the indicator of the all-plus vertex.
Consequently
\[
 \widehat f_n(S)=
 \begin{cases}
  1/2+\varepsilon,&S=\varnothing,\\
  -1/2+\varepsilon,&\varnothing\ne S\subseteq\{1,2\},\\
  \varepsilon,&S\not\subseteq\{1,2\}.
 \end{cases}                                             \tag{4}
\]
In particular, \(\widehat f_n([n])=\varepsilon\ne0\), so
\(\deg f_n=n\).

Formula (4) gives immediately
\[
 \boxed{H_n=L(q_+)+3L(q_-)+(2^n-4)L(r)},                 \tag{5}
\]
and
\[
 \boxed{I_n=1+(n-4)\varepsilon}.                        \tag{6}
\]

\subsection*{Exact upper-set functional}

Write a Fourier label as \(S=A\cup B\), where
\(A\subseteq\{1,2\}\) and
\(B\subseteq\{3,\ldots,n\}\).  If \(B\ne\varnothing\), every
coefficient above \(S\) has mass \(r\), hence
\[
 \pi(S)=2^{n-|S|}r.
\]
For the two exceptional types one has
\[
 \pi(\{1\})=\pi(\{2\})=\frac12-\varepsilon,
 \qquad
 \pi(\{1,2\})=\frac14-\frac\varepsilon2.
\]
Summing every label, including the empty and terminal rows, yields
\[
 \boxed{
 U_n=L(q_+)+L(q_-)
      +q_-\log_2\frac1{1-2\varepsilon}
      +n\varepsilon-(4n-4)\varepsilon^2.}               \tag{7}
\]
No normalization or zero-mass term is omitted.

\subsection*{Exact restriction aggregate}

Let \(V_N\) denote \(U\) for a dictator on \(N\) variables whose
all-plus value has been flipped.  With \(e=2^{1-N}\), the same upper-set
calculation gives
\[
 V_N=L(e^2)+\frac12L((1-e)^2)+Ne-(2N-1)e^2.              \tag{8}
\]
Restricting either \(x_1\) or \(x_2\) gives one constant child and one
such punctured dictator on \(n-1\) variables.  Restricting any of the
other \(n-2\) coordinates gives one unpunctured two-bit defect, whose
\(U\)-value is exactly one, and one copy of \(f_{n-1}\).  Therefore, for
\(n\ge4\),
\[
 \boxed{
 G_n=nU_n-V_{n-1}-\frac{n-2}{2}(1+U_{n-1}).}             \tag{9}
\]
The endpoint \(n=3\) has a constant punctured codimension-two child and is
handled directly; (5)--(8) remain valid and the checker records its exact
value.

\subsection*{High-degree reserve}

Here
\[
 M_1=\frac12-2\varepsilon+n\varepsilon^2,
 \qquad
 W_2=\frac14-\varepsilon+\binom n2\varepsilon^2.
\]
Hence
\[
 \boxed{
 Q_{\ge3}=\frac14+(n-1)\varepsilon
       -\left(n+\binom n2\right)\varepsilon^2.}          \tag{10}
\]
This retains \(W_1-M_1\), \(W_2\), and every degree-at-least-three
contribution exactly.

\subsection*{Tensor primeness}

For any nontrivial variable partition \(A\sqcup B=[n]\), a factorization
\(f_n=g(x_A)h(x_B)\) would make the Fourier coefficient matrix across
\(A\mid B\) have rank one.  This never occurs.  If auxiliary variables
\(i\in A\cap\{3,\ldots,n\}\) and
\(j\in B\cap\{3,\ldots,n\}\) exist, the minor indexed by
\(\varnothing,\{i\}\) and \(\varnothing,\{j\}\) has determinant
\[
 \varepsilon(1/2+\varepsilon)-\varepsilon^2
 =\varepsilon/2\ne0.
\]
If all auxiliary variables lie on one side, choose an auxiliary variable
there and one of \(x_1,x_2\) on the other side.  The analogous determinant
is
\[
 \varepsilon[(1/2+\varepsilon)-(-1/2+\varepsilon)]
 =\varepsilon\ne0.
\]
Thus \(f_n\) is tensor-prime for every \(n\ge3\).

\subsection*{Controlled asymptotics and verdict}

Taylor expansion of the displayed exact formulas gives, uniformly as
\(n\to\infty\),
\[
 \begin{aligned}
 H_n&=2+\left(4n-8+\frac2{\ln2}\right)\varepsilon
          +O(n\varepsilon^2),\\
 U_n&=1+\left(n+\frac1{2\ln2}\right)\varepsilon
          +O(n\varepsilon^2),\\
 G_n&=2+\left(n-\frac1{\ln2}\right)\varepsilon
          +O(n^2\varepsilon^2),\\
 I_n&=1+(n-4)\varepsilon.
 \end{aligned}                                           \tag{11}
\]
Consequently
\[
 \frac{H_n}{I_n}\to2,\qquad
 \frac{U_n}{I_n}\to1,\qquad
 \frac{G_n}{I_n}\to2,\qquad
 \frac{Q_{\ge3}(f_n)}{I_n}\to\frac14.                  \tag{12}
\]
For the frozen pair,
\[
 G_n-4I_n-\frac34U_n
 =-\frac{11}{4}+O(n\varepsilon),                         \tag{13}
\]
so this family is not a GU falsifier and gives no unbounded
\(G/I\), \(U/I\), or \(H/I\) sequence.

The structural information is instead sharp: tensor-prime functions of
arbitrarily high degree can converge spectrally to a degree-two Boolean
boundary while \(Q_{\ge3}/I\to1/4\).  Therefore a finite degree-at-least-
three census cannot export its larger strict reserve margin unchanged to
all dimensions.  Any reserve-only tail theorem must allow a constant at
most \(1/4\).  The family neither disproves the possible universal bound
\(Q_{\ge3}\ge I/4\) nor proves the signed leaf (3); those are the precise
remaining obligations.
